\chapter{Background}\label{ch:background}
% =============================================================================

This chapter introduces the concepts behind each contribution of \autoref{sec:intro:contributions}.
Their mathematical form is given in \autoref{ch:method}.

\section{Learning Trajectory Policies from Demonstration}\label{sec:bg:il}

The setting is a dataset of expert trajectories, a conditional model over horizon-length
trajectory segments \parencite{janner2022planning,chi2023diffusion}, and receding-horizon
execution: the model plans a short horizon, the first action is executed, and the model plans again.
The benchmarks come from \parencite{jia2024towards}, whose demonstrations are multimodal. Planning
over a horizon is what makes constraint enforcement possible: a single action has nothing to project.
Multimodality is why the model is generative and not a regressor, and it is also why a truncated
sampler risks averaging between modes.

\section{Diffusion Models and Diffusion Predictive Control}\label{sec:bg:diffusion}

A denoising diffusion model \parencite{sohl2015deep,ho2020denoising} learns a distribution by
reversing a fixed process that gradually adds Gaussian noise to the data. A network is trained to
predict the noise added at each step of a noise schedule; generation starts from pure noise and
removes it step by step. Diffuser \parencite{janner2022planning} applies this to whole trajectories,
and diffusion policies \parencite{chi2023diffusion} brought it to visuomotor control.

Two properties follow from the construction. The number of denoising steps is tied to the schedule
the network was trained on, so it is fixed when training ends; fewer steps require a different sampler
\parencite{song2021denoising} or a distilled model. And each step adds fresh noise, so two plans drawn
for the same observation differ.

\ac{DPCC} \parencite{romer2025diffusion} makes such a planner respect state and action constraints
that are absent from the demonstrations. Between denoising steps it projects the intermediate
trajectory onto the set of trajectories that satisfy the constraints and a model of the dynamics, and
justifies the projection as guidance towards that set. It shrinks the constraint sets by a margin that
absorbs the mismatch between the dynamics model and the real system, and it draws several candidate
plans per control step and selects one.

\section{Flow Matching}\label{sec:bg:fm}

Flow matching \parencite{lipman2023flow,lipman2024guide} describes generation as transport: a
time-dependent velocity field moves samples from a noise distribution to the data distribution, and a
sample is produced by integrating an ordinary differential equation along that field. The field is
learned by regression. Because the field of the whole distribution is intractable, training regresses
the velocities of simple per-sample paths between a noise draw and a data point, which are known in
closed form and yield the same gradient \parencite{lipman2023flow}. With a straight-line path, as in
rectified flow \parencite{liu2023flow}, the regression target is the difference between the data point
and the noise, and no trajectory has to be simulated during training. Flow matching has achieved
state-of-the-art results across images, video, audio, speech and biological structures
\parencite{lipman2024guide}.

Two consequences matter for control. The number of integration steps is a property of the solver, not
of training, so it is chosen at deployment. And the integration adds no noise: a plan is a
deterministic function of its initial noise and the observation.

\section{Few-Step Generative Models}\label{sec:bg:fewstep}

A step budget chosen at deployment is not automatically a good one. Straight per-sample paths average
into a curved field, and coarse integration steps along a curved field are inaccurate however well the
network is trained \parencite[Sec.~3]{geng2025meanflow}. Few-step models change what the network
learns so that a large step stays accurate.

MeanFlow \parencite{geng2025meanflow} learns the average velocity over a time interval instead of the
instantaneous velocity. One evaluation of an accurate average velocity covers the whole interval, so a
plan can be generated in a single step. The average velocity is defined by an integral that cannot be
evaluated during training; MeanFlow differentiates that definition into an identity relating the
average velocity to the instantaneous velocity and to its own time derivative, and computes the
derivative with a \ac{JVP}. When the interval shrinks to a point, the MeanFlow target becomes the
flow-matching target.

A subsequent analysis \parencite{zhang2025alphaflow} decomposes the MeanFlow objective into a
trajectory flow-matching term and a trajectory-consistency term, and builds a family of objectives
that anneals from the first to the second. The target of this family is formed at an intermediate time
from a stop-gradient copy of the network itself, the construction of consistency models
\parencite{song2023consistency}; this thesis refers to it as \emph{consistency training}. MeanFlow and
this family of objectives both achieve state-of-the-art one- and two-step image generation
\parencite{geng2025meanflow,zhang2025alphaflow}.

For diffusion the step budget is therefore fixed at training time, for flow matching it is chosen at
deployment, and MeanFlow and consistency training make a small budget accurate.

\section{Constraint Enforcement in Generative Sampling}\label{sec:bg:mpc}

Model predictive control \parencite{morari1999model} optimises over a finite horizon, executes the
first action and optimises again. A learned planner fits this loop, but its plans carry no guarantee
for constraints the demonstrations never contained. Three approaches exist. Guidance adds the gradient
of a constraint cost to the sampler \parencite{carvalho2025motion}, which makes violations less likely
without excluding them. Projection maps a sample onto the feasible set, after sampling
\parencite{power2023sampling} or at every step \parencite{christopher2024constrained,romer2025diffusion};
the projection of \ac{DPCC} includes a model of the dynamics, so the projected plan stays dynamically
feasible. Optimisation-based methods steer the sampler itself.

One such method \parencite{li2025hardflow} steers a flow-matching sampler by treating sampling as an
optimal control problem. A control input is added to the learned velocity field; its effort is penalised, so
the controlled sampler stays close to the learned one; and the constraints are imposed on the final
sample only, not on the path that leads to it. Solving that problem exactly would require simulating
the rest of the trajectory at every step. The method avoids this in three moves. It splits the problem,
as model predictive control does, into one short problem per sampling step. It replaces the unknown
final sample by its expected value given the current state, which for a flow-matching model follows in
closed form from the state and the learned velocity. And it optimises directly over that predicted
final sample instead of over the next state, so the constraints act on a plan and not on its image
through the network. The corrected prediction is then mapped back into the next state of the sampler.
At the final step the prediction and the state coincide, so the final sample satisfies the constraints
exactly. An optional terminal cost can further shape the sample, and the optimisation is applied only
in the later sampling steps, where the prediction of the final sample is accurate.

The two approaches therefore differ in where the constraints act. Per-step projection corrects the
sampler's current iterate, which early in sampling is still far from any plan; this thesis calls it
\emph{iterate projection}. The optimal-control formulation corrects the plan the sampler is heading
towards; this thesis calls it \emph{endpoint projection}. Plans generated with neither are \emph{unguided}.

\section{Vision-Conditioned Policies}\label{sec:bg:visual}

A policy that observes images compresses each camera view into a feature vector with a convolutional
encoder and conditions the generative model on that vector. D3IL \parencite{jia2024towards} provides
such a vision pipeline for its tasks, in which ResNet encoders process each camera view. Concatenated
with the time embedding of the network, the features reach its layers as a feature-wise conditional
bias \parencite[Sec.~3]{perez2017film}. Transformer backbones operate on image patches instead of a
pooled feature vector \parencite{peebles2022dit}.

\section{Quadrotor Dynamics and Control}\label{sec:bg:quadrotor}

A manipulator in position control is fully actuated and holds its configuration without effort; a
planned sequence of end-effector positions becomes joint commands through \ac{IK}. A quadrotor is
different. Its four rotors produce a thrust along one body axis and three moments, so it controls four
of its six degrees of freedom directly and must tilt in order to translate, and hover is an unstable
equilibrium that requires continuous feedback. Such an underactuated, open-loop unstable vehicle is
controlled in cascade: an outer loop turns a position error into a desired thrust vector, which fixes a
desired attitude, and an inner loop tracks that attitude. Geometric controllers formulate the attitude
error directly on the rotation group and thereby avoid the singularities of local coordinates
\parencite{lee2010geometric}. Model predictive control is an alternative that optimises the rotor
commands over a short horizon against a simulation of the vehicle; predictive sampling in \ac{MJPC}
does so without derivatives, by rolling out perturbed control sequences and keeping the best
\parencite{howell2022predictive}. The vehicle and its scenes are simulated in MuJoCo
\parencite{todorov2012mujoco}.

% =============================================================================
